% \documentclass[12pt,twoside]{report}
\documentclass[12pt]{report}
\usepackage[utf8]{inputenc}
\usepackage{graphicx}
\graphicspath{{images/}}
\usepackage[a4paper,width=150mm,top=25mm,bottom=25mm,bindingoffset=6mm]{geometry}
\usepackage{setspace}
\usepackage{fancyhdr}
\usepackage{listings}
\usepackage{parcolumns}
\usepackage{color}
\usepackage{float}

\usepackage[bottom]{footmisc}

\usepackage{enumitem}

\newlist{questions}{enumerate}{2}
\setlist[questions,1]{label=RQ\arabic*.,ref=RQ\arabic*}
\setlist[questions,2]{label=(\alph*),ref=\thequestionsi(\alph*)}

\usepackage{hyperref}
\usepackage[capitalise]{cleveref}
\hypersetup{
    colorlinks=true,
    linkcolor=blue,
    filecolor=magenta,      
    urlcolor=cyan,
    pdftitle={Overleaf Example},
    pdfpagemode=FullScreen,
    }

\definecolor{codegreen}{rgb}{0,0.6,0}
\definecolor{codegray}{rgb}{0.5,0.5,0.5}
\definecolor{codepurple}{rgb}{0.58,0,0.82}
\definecolor{backcolour}{rgb}{0.95,0.95,0.92}

\lstdefinestyle{mystyle}{
    backgroundcolor=\color{backcolour},   
    commentstyle=\color{codegreen},
    keywordstyle=\color{magenta},
    numberstyle=\tiny\color{codegray},
    stringstyle=\color{codepurple},
    basicstyle=\ttfamily\footnotesize,
    breakatwhitespace=false,         
    breaklines=true,     
    postbreak=\mbox{\textcolor{red}{$\hookrightarrow$}\space},            
    captionpos=b,                    
    keepspaces=true,                 
    numbers=left,                    
    numbersep=5pt,                  
    showspaces=false,                
    showstringspaces=false,
    showtabs=false,                  
    tabsize=2
}

\usepackage{color}
\definecolor{gray}{rgb}{0.4,0.4,0.4}
\definecolor{darkblue}{rgb}{0.0,0.0,0.6}
\definecolor{cyan}{rgb}{0.0,0.6,0.6}

\lstset{
  basicstyle=\ttfamily,
  columns=fullflexible,
  showstringspaces=false,
  commentstyle=\color{gray}\upshape
}

\lstdefinelanguage{XML}
{
  morestring=[b]",
  morestring=[s]{>}{<},
  morecomment=[s]{<?}{?>},
  stringstyle=\color{black},
  identifierstyle=\color{darkblue},
  keywordstyle=\color{cyan},
  morekeywords={xmlns,version,type}% list your attributes here
  breaklines=true,     
    postbreak=\mbox{\textcolor{red}{$\hookrightarrow$}\space},  
}

\usepackage{chngcntr}
\counterwithout{figure}{chapter}
\counterwithout{table}{chapter}

\usepackage[activate={true,nocompatibility},final,tracking=true,kerning=true,spacing=true,factor=1100,stretch=10,shrink=10]{microtype}

\usepackage[backend=biber,sorting=none]{biblatex}
\addbibresource{references.bib}

\title{
    {Inferring security-relevant properties from Java documentation}\\
    {\large University of Queensland}\\
    {\includegraphics{images/uq_logo.jpg}}
}

\author{Connor Myers}
\date{\today}

%\doublespacing
\onehalfspacing

\begin{document}

% \maketitle
\begin{titlepage}
    \begin{center}
        \vspace*{1cm}

        \includegraphics[width=0.5\textwidth]{uq_logo.jpg}

        \vspace{2.5cm}  

        \LARGE
        \textbf{Inferring Security-relevant Properties from Java Documentation}

        \vspace{0.5cm}

        \large
        \textbf{Connor Myers} \\
%        \vspace{0.25cm}
%        \normalsize
%        45323979


        \vspace{0.75cm}
        \normalsize
        Supervised by\\
        \vspace{0.25cm}
        \large
        Mark Utting, Behnaz Hassanshahi, and Dai Nguyen


        % \vspace{3.5cm}


        % \vspace{1cm}

        \vfill

        School of Information Technology and Electrical Engineering University of Queensland\\

        \vspace{1.5cm}

        Submitted for the degree of Bachelor of Engineering (Honours)\\ in the division of Software Engineering

        \vspace{0.8cm}

        \monthyeardate\today

    \end{center}
\end{titlepage}
\thispagestyle{empty}

\begin{flushright}
    Connor Myers \\
    connor.myers@uqconnect.edu.au \\ 
    45323979
\end{flushright}

\vspace{2cm}

\noindent
\today

\vspace{1cm}

\noindent
Prof Michael Brünig, \\
\noindent
Acting Head of School \\
\noindent
School of Information Technology and Electrical Engineering \\ 
\noindent
The University of Queensland \\
\noindent
St Lucia  QLD  4072

\vspace{1cm}

\noindent
Dear Professor Brünig, \\

\noindent
In accordance with the requirements of the Degree of {Bachelor of Engineering (Honours)} in the School of Information Technology and Electrical Engineering, I submit the following thesis entitled ``Inferring Security Relevant Properties from Java Documentation''. \\

\noindent
The thesis was performed under the supervision of Associate Professor Mark Utting. I declare that the work submitted in the thesis is my own, except as acknowledged in the text and footnotes, and that it has not previously been submitted for a degree at the University of Queensland or any other institution. \\

\noindent
Yours sincerely \\

\noindent
Connor Myers

\pagenumbering{Roman}
\chapter*{Abstract}
\setcounter{page}{1}

\noindent
Javadoc, while being an informal natural language description of Java code, may contain information relevant for code security tasks. In this thesis natural language processing techniques are applied on Javadoc to automatically classify security-sensitive Java methods for the purpose of configuring taint analysis tools. \\

\noindent
A variety of natural-language processing techniques and machine learning models were explored. In particular, pre-trained bimodal machine learning models for understanding code and natural language (CodeBERT and UniXcoder) were fine-tuned on a dataset of labelled security-sensitive methods. The performances of these machine learning models were evaluated on a realistic test dataset. \\

\noindent
The training dataset produced for this thesis was comprised of 931 labelled security-relevant methods. The security-relevant labels used for this dataset  were: sources (125), sinks (230), sanitisers (156), and security-irrelevant methods (420). The dataset is publicly available \href{https://github.com/uqcyber/SecurityMethodClassifier/tree/main/tools/dataset-manager}{here}. \\

\noindent
In addition to the training dataset, a test dataset was produced to more accurately reflect a realistic use-case for these models. The dataset was produced by manually classifying a random sample of Java methods from Apache's Log4J. \\


\noindent 
The highest performing model configuration achieved an F1-score of $0.68$ on the test dataset. This exceeded the baseline performance of the naive approach to this task which had an F1-score of $0.51$.

%\newpage
%
%\vspace*{0.5\textheight}
%
%\begin{center}
%\emph{For Jeff}
%\end{center}

\chapter*{Acknowledgements}
I would like to thank my three supervisors Mark Utting, Behnaz Hassanshahi, and Dai Nguyen for going above and beyond with helping me with thesis. I would also like to thank Paddy Krishnan from Oracle for organising this collaboration between UQ and Oracle.

\noindent

\tableofcontents

\cleardoublepage\pagenumbering{arabic}
\chapter{Introduction}
\noindent
Static application security testing (SAST) applies static program analysis to the security of software projects. Static analysis is performed on the source code, compiled byte-code, or some other representation of the source code \cite{principles_program_analysis}. Popular SAST tools include: SonarQube, SpotBugs, Snyk Code, and Facebook's Infer. SAST tools have become widely used at all stages of the software development life-cycle. Furthermore, they have been demonstrated to be highly effective in identifying defects in code compared to other methods, such as manual inspection \cite{costs_static_analysis}.\\

\noindent
Unfortunately, SAST tools are not perfect. From Rice's Theorem, it is known that any non-trivial property about programs is reducible to the halting problem, i.e., undecidable \cite{moller2012static}. That is, a perfect analysis is not always possible, and approximations must be made in order for the analysis to terminate, resulting in an imperfect output. False positives (where safe code is incorrectly labelled as a defect) and false negatives (where the analysis misses a genuine fault) \cite{chess2004static} are common in all static analysis tools. False positives are not ideal as manual inspection is required to isolate genuine defects, leading to an increased time and labour cost. And, false negatives of course mean that genuine software faults are left undetected by the analysis. While solving the halting problem may be out of scope for an undergraduate thesis, there are still additional methodologies for improving existing static analyses that do not involve upsetting Mr. Turing.\\

\noindent
To reduce the rate of false negatives, manual configuration of the SAST tools for specific software projects is often required \cite{ZampettiFiorella2017HOSP}. For example, to perform an effective taint analysis, a custom configuration of sources, sinks and sanitisers for a particular project is often required. These taint analysis tools are not smart enough to know if a method is a source, sink, or a sanitiser --- they require sensitive methods to be pre-labelled to conduct the analysis. Generally, these tools come with a default configuration containing a list of sensitive methods in popular APIs. For example, the getParamater() method as a part of the ServletRequest interface would likely be listed by all taint analysis tools as a source. Of course, if a project is using a niche library or code written in-house, the taint analysis will not be able to incorporate those data-points into its analysis, leading to a sub-optimal analysis. To improve the taint analysis, one can always identify the sensitive methods manually and add these to the taint analysises configuration. However, creating this manual configuration is often a time-consuming and resource-intensive process. As such, methodologies to automatically label security sensitive methods may enable existing taint analysis tools to be configured more efficiently and effectively.\\

\noindent
Deep learning, a subtopic within the broader machine learning, has recently increased in popularity. What was once a niche topic reserved for academics with access to computational resources far out of reach by most has become arguably one of the most influential technologies of the modern age. Deep learning has fundamentally changed our everyday life, with applications ranging from targeted advertising to identifying COVID-19 cases from lung CT scans \cite{RAHIMZADEH2021102588}. Natural-language processing (NLP) is one often overlooked application of deep learning, but it has seen tremendous growth in recent years.\\

\chapter{Background Knowledge}
Lorem ipsum dolor sit amet, consectetur adipiscing elit. Vivamus nulla libero, semper id dui finibus, suscipit pretium purus. Duis vehicula eleifend justo, at cursus turpis vulputate eget. Nullam vestibulum id neque a convallis. Sed molestie vehicula turpis, vel aliquet felis efficitur a. Fusce accumsan metus augue, id sagittis eros elementum vitae. Cras vel ultrices odio. Sed ac augue ex. Etiam tristique dignissim nibh, sit amet tincidunt libero tempor sit amet. Suspendisse tincidunt vitae ligula a posuere. Morbi euismod felis quis risus malesuada cursus. Curabitur imperdiet pretium odio, ac posuere turpis vehicula in. Praesent venenatis urna in augue congue elementum. Nam ligula nibh, rhoncus vitae luctus vel, feugiat sed neque. Curabitur ut arcu sed magna vulputate sollicitudin id a libero. Praesent tempus leo vel eros ornare mollis. Sed suscipit at mi ut viverra.

Maecenas vitae tempus arcu. Nullam maximus id ligula a finibus. Aliquam laoreet faucibus sem. Mauris interdum erat vel pretium iaculis. Sed et pretium arcu, at condimentum lorem. Suspendisse tincidunt neque vel ullamcorper consequat. Phasellus eget ligula quis velit fringilla porttitor. Pellentesque quis risus id erat euismod bibendum mattis nec sapien. Duis gravida, turpis at dapibus consequat, ex mi tempor augue, in scelerisque mi nisi ut lorem. Donec orci tellus, lobortis ac pretium non, elementum et dolor. Morbi sed hendrerit erat. Praesent nec ante lorem. Vestibulum dapibus varius nisl at hendrerit. Praesent aliquam vehicula dolor et pulvinar. Aenean facilisis leo eget lacus rhoncus, ac tempus urna molestie. Vivamus efficitur erat quis felis pulvinar mattis quis volutpat lorem.

Nullam luctus velit varius massa suscipit fringilla. Quisque egestas diam nec eros sagittis posuere. Sed venenatis facilisis nibh ac mattis. Morbi a mi at est vehicula imperdiet. Fusce tempus justo at tortor volutpat vulputate. Vestibulum molestie congue interdum. Nullam auctor faucibus risus ac hendrerit. Nulla viverra fermentum vulputate. Nulla mattis elementum ligula. Fusce eu tellus sapien. Maecenas scelerisque odio lacus, quis venenatis sapien molestie semper. Donec in consequat libero, quis euismod tellus.

Vestibulum eget lobortis lorem, id aliquam leo. Morbi ornare, enim et iaculis luctus, erat dui rutrum ligula, vel aliquam libero nisi vitae eros. Proin sodales ipsum nec ipsum tristique, porttitor sagittis tortor rhoncus. Integer auctor rutrum imperdiet. Nulla sollicitudin facilisis mauris, in mattis lorem ultrices quis. Curabitur neque lectus, auctor at ullamcorper et, porta eu nulla. Cras efficitur felis in ipsum volutpat, vitae condimentum velit ornare.

Donec consectetur, libero tempus sollicitudin gravida, felis erat aliquam libero, viverra maximus lacus nisi in lacus. Aenean hendrerit volutpat felis, sed viverra ex accumsan a. Vivamus vel tortor viverra nisi tincidunt dapibus. Proin diam libero, varius at turpis et, tincidunt congue arcu. Nunc tristique elit libero, nec rutrum quam aliquet at. Orci varius natoque penatibus et magnis dis parturient montes, nascetur ridiculus mus. Vestibulum at nisi quis nulla malesuada tempor. Mauris eu felis risus. Vestibulum ante ipsum primis in faucibus orci luctus et ultrices posuere cubilia curae; Donec convallis est ex, vitae facilisis eros commodo in. Nullam faucibus nunc ac magna congue, et porttitor magna eleifend. Nunc at velit ullamcorper, pellentesque lorem in, dictum arcu. Suspendisse nec auctor velit. Mauris congue magna molestie augue maximus, ac hendrerit ante viverra. Pellentesque convallis nisl et felis venenatis varius.

\chapter{Literature Review}
\section{Similar Tasks}
There are numerous similar tasks in the literature. That is, applying machine-learning to the classification of security-sensitive methods. Here we will detail the more relevant papers in the literature.

\subsection{Swan}

SWAN is an automated classifier of security-sensitive methods in Java using machine-learning \cite{swanassist}. SWAN classifies code as a source, sink, sanitizer, authentication-method, or none --- which is identical to our task, with the inclusion of the \emph{authenitication-method} security-sensitive method type. SWAN is also able to further classify some of these methods according to their specific CWE (Common Weakness Enumeration), however this aspect of SWAN is not relevant to this thesis. \\

 \noindent
 SWAN employs a Support Vector Machine (SVM) approach on 206 manually defined binary features (i.e., does the method have this feature --- true or false?) to perform its classification. That is, the authors have gone through a dataset of examples of security-sensitive-methods and manually identified what aspects of a code is important for classifying it as a security sensitive method. This is decidedly different from our proposed implementation of fine-tuning a pre-trained bimodal Transformer model. \\

\noindent
For example, one of the 206 features SWAN classifies on is whether or not the class name the method belongs to contains the keyword `SQL' (the reason for this being that methods interacting with an SQL database are more likely to be a \emph{sink}). While this may seem like a simple feature, which it is, when combined with 205 additional features a high accuracy can be reached on a relatively small training dataset.\\

\noindent
SWAN's training dataset is publicly available and open-source \cite{swan_github}. Labelled datasets of security-sensitive methods are hard to come by, and thankfully Swan has quite an extensive list. For this thesis, we will be including a subset of this dataset (we remove \emph{authentication-method} examples, add \emph{none} examples which are not included in SWAN's public dataset, and remove unsuitable examples according to our own filtering rules) for our own training and evaluation purposes. Furthermore, SWAN does not provide its testing or validation dataset so we will split up their training set to get testing and validation examples.\\

\noindent
This dataset does not provide the code and Javadoc associated with each method, so additional processing of this data is required. We will come back to this later in the thesis. \\

\noindent
The \emph{SWAN} project also includes SWAN-assist, which is an Intellij plugin that allows users to classify security-relevant methods while working on a project. SWAN includes this additional information to improve the machine-learning classifier for a particular project. Of course, this aspect is not relevant to our project but it does inflate the accuracy of the model when compared to a fully-automated process such as ours. \\

\noindent
SWAN-base (i.e., SWAN without additional information from users using SWAN-assist) achieves an accuracy of $0.826$ on classification of security-sensitive methods \cite{swanassist}. While it would be good for our model to achieve an accuracy higher than this, it should be noted that it is not a fair comparison. For one, we will be operating on a subset of SWAN's base training-set out of necessity, and will be classifying on four labels, not five. Naturally, classifying on fewer labels is likely to inflate the accuracy of our model compared to SWAN-base, whereas operating on a subset on the dataset is likely to reduce the accuracy of our model relatively.

\subsection{SuSi}



\chapter{Collecting the Data}
Lorem ipsum dolor sit amet, consectetur adipiscing elit. Vivamus nulla libero, semper id dui finibus, suscipit pretium purus. Duis vehicula eleifend justo, at cursus turpis vulputate eget. Nullam vestibulum id neque a convallis. Sed molestie vehicula turpis, vel aliquet felis efficitur a. Fusce accumsan metus augue, id sagittis eros elementum vitae. Cras vel ultrices odio. Sed ac augue ex. Etiam tristique dignissim nibh, sit amet tincidunt libero tempor sit amet. Suspendisse tincidunt vitae ligula a posuere. Morbi euismod felis quis risus malesuada cursus. Curabitur imperdiet pretium odio, ac posuere turpis vehicula in. Praesent venenatis urna in augue congue elementum. Nam ligula nibh, rhoncus vitae luctus vel, feugiat sed neque. Curabitur ut arcu sed magna vulputate sollicitudin id a libero. Praesent tempus leo vel eros ornare mollis. Sed suscipit at mi ut viverra.


\chapter{Models}
\section{Model Architecture Summary}

In Section \ref{model_architecture} a variety of machine-learning architectures and techniques were discussed. In this section, it will be made explicit the exact machine-learning architectures that will be evaluated. There will be four main categories made up of different variations.

\subsection{SVM with TF-IDF}

\noindent
In this model, features will be generated using the the primitive text-frequency inverse-document-frequency methodology. Features generated by tf-idf will be fed into a linear Support Vector Machine classification head. \\

\noindent 
The sci-kit learn Python library will be used both for the tf-idf and SVM classifier functionality. \\

\noindent
This model was chosen as one of our architectures as it represents arguably the most simplistic way this problem could be approached. In other words, the performance of this model should be treated as the `base-line' performance for which all of the additional more-sophisticated models should at least reach. 

\subsection{Pre-trained SVM-Transformer Hybrid}

\noindent
In this model, features will be extracted using CodeBERT or UniXcoder (both configurations will be tested) in there base pre-trained form (\emph{no} additional fine-tuning). Sentence embeddings will be pooled from the contextualised token embeddings (mean and CLS will be tested). An SVM classifier will then classify on these sentence embeddings. \\

\noindent
To load CodeBERT and UniXcoder, as well as to extract the features and pool the token embeddings, the machine-learning framework PyTorch will be used. As with other models, the linear SVM classifier functitonality will come from sci-kit learn. \\

\noindent
The purpose of this model is two-fold. Firstly, it would be good to know how CodeBERT and UniXcoder perform with no fine-tuning. The results can be compared to the performance of the model when it is fine-tuned, and from that an estimate to how much performance was gained from fine-tuning can be established. Secondly, due to the dataset being quite small, it is possible there are too few training examples to properly train a neural-network classification head. Recall that only the \emph{encoder} has been pre-trained, whereas the classification-head parameters will be trained from scratch on our dataset alone. With this constraint, it is possible that a more simple, linear classification head may perform better (i.e. a SVM classifier).

\subsection{Fine-tuned SVM-Transformer Hybrid}

\noindent
This architecture is identical to the previous pre-trained SVM-transofrmer hybrid, except the encoders will first be fine-tuned on our dataset (using a neural-network classification-head) and then a linear SVM classifier will be trained with features extracted from the training dataset using the fine-tuned encoders. To clairfy: the training dataset for the SVM classifier and encoders are identical. \\

\noindent
As with the similar model using pre-trained encoders, PyTorch will be used the encoders, and sci-kit learn will be used for the SVM classifier.\\

\noindent
This approach will be tested as it is believed fine-tuning the encoder first will enable the encoders to generate better embeddings for the SVM classifier. Furthermore, because the dataset is small, a linear SVM classifier may perform better than a neural-network classifier.\\

\subsection{Fine-tuned Transformer with neural-network classifier}

\noindent
This architecture is most reflective of how CodeBERT and UniXcoder were intended to be used. Pooled sentence embeddings created by the encoders are fed into a neural-network classification head, which generates a probability distribution for the likelihood of the input data being one of th 4 labels. The label with the highest priority is selected.\\

\noindent
PyTorch will be used to implement this model.\\

\noindent
This approach will be tested as it is the most sophisticated of all 4 models that will be tested. It is believed this model will perform the best. 

\subsection{Training Settings}

This section details the important hyperparameters used for the machine-learning models.

\subsubsection{SVM Classifier}

\begin{itemize}
    \item \textbf{Loss function}: hinge loss
    \item \textbf{Penalty:} L2 Regularisation
    \item \textbf{Lagrangian Multiplier:} $e^{-3}$
    \item \textbf{No. Epochs:} $5$
\end{itemize}

\subsubsection{CodeBERT and UniXcoder}

\begin{itemize}
    \item \footnote{This low training batch-size is due to VRAM limitations of the testing machine.}\textbf{Training batch size:} $4$
    \item \textbf{Validation batch size:} $32$
    \item \textbf{Learning Rate: } $5e^{-5}$
    \item \textbf{Max gradient norm: } $1.0$
    \item \footnote{It was experimentally determined that no extra accuracy or decrease in loss was obtained for any additional iterations}\textbf{No. Epochs:} $3$
\end{itemize}

\section{Performance on Validation Dataset}

In this section, the performance of the \emph{trained} models will be evaluated on the validation dataset. As established in Section \ref{dataset_limitations}, these results will be inflated compared to a real-world use-case. \\

\noindent
The various tables below show the performance of each model (and their multiple cofigurations, if applicable) on the validation dataset.

% \subsection{SVM with tf-idf}

\begin{table}[H]
    \begin{center}
        \begin{tabular}{c|c|c}
            \textbf{Feature Generator} & \textbf{Classifier} & \textbf{Accuracy} \\
            \hline
            TF-IDF            & SVM        & 0.66    
        \end{tabular}
    \end{center}
    \caption{Validation performance of SVM-classifier with tf-idf features model}
    \label{tab:validation_perfomance_svm_tfidf}
\end{table}

% \subsection{Pre-trained SVM-Transformer Hybrid}

\begin{table}[H]
    \begin{center}
        \begin{tabular}{c|c|c}
            \textbf{Feature Generator} & \textbf{Pooling Type} & \textbf{Accuracy }\\
            \hline
            CodeBERT & CLS & 0.27 \\
            CodeBERT & Mean & 0.39 \\
            UniXcoder & CLS & 0.78 \\
            UniXcoder & Mean & 0.76 \\
        \end{tabular}
    \end{center}
    \caption{Validation performance of SVM-classifier with tf-idf features model}
    \label{tab:validation_perfomance_pre_svm}
\end{table}

% \subsection{Fine-tuned SVM-Transformer Hybrid}

\begin{table}[H]
    \begin{center}
        \begin{tabular}{c|c|c}
            \textbf{Feature Generator} & \textbf{Pooling Type} & \textbf{Accuracy }\\
            \hline
            CodeBERT & CLS & 0.79 \\
            CodeBERT & Mean & 0.79 \\
            UniXcoder & CLS & 0.79 \\
            UniXcoder & Mean & 0.78 \\
        \end{tabular}
    \end{center}
    \caption{Validation performance of SVM-classifier with tf-idf features model}
    \label{tab:validation_perfomance_fine_svm}
\end{table}

% \subsection{Fine-tuned Transformer with neural-network classifier}

\begin{table}[H]
    \begin{center}
        \begin{tabular}{c|c|c|c}
            \textbf{Feature Generator} & \textbf{Pooling Type} & \textbf{Accuracy} & \textbf{Loss} \\
            \hline
            CodeBERT & CLS & 0.79 & 1.08\\
            CodeBERT & Mean & 0.77 & 1.49\\
            UniXcoder & CLS & 0.80 & 1.08\\
            UniXcoder & Mean & 0.79 & 1.26\\
        \end{tabular}
    \end{center}
    \caption{Validation performance of SVM-classifier with tf-idf features model}
    \label{tab:validation_perfomance_fine_nn}
\end{table}



\chapter{Experiments}
\section{Performance on Realistic Use-case}

\subsection{Experiment Description}

As mentioned in Section \ref{dataset_limitations}, the performance of the models on the validation dataset is not exactly reflective of a realistic use case.\\

\noindent
To alleviate this concern, we conduct a test experiment on a realistic use case. That is, we test all model configurations on the test dataset that was described in Section \ref{test-dataset}. \\

%The test dataset can be found at \url{www.google.com}.\\

\noindent
For these test experiments, F1-score\footnote{F1-score is the harmonic mean of precision and recall} will be used rather than accuracy as the primary metric to gauge performance.

\subsection{Experiment Results}
\label{experiment_results}

\noindent
\crefrange{tab:test_perfomance_tfidf_svm}{tab:test_perfomance_fine_nn} show the performance of the various model configurations on the Log4J test dataset. \\

\begin{table}[H]
\caption{Performance of SVM-classifier with tf-idf feature generation on the test dataset}
\label{tab:test_perfomance_tfidf_svm}
    \begin{center}
        \begin{tabular}{c|c|c|c}
            \textbf{Accuracy} & \textbf{Precision} & \textbf{Recall} & \textbf{F1-score}  \\
            \hline
            0.50 & 0.53 & 0.50 & 0.51\\
        \end{tabular}
    \end{center}
\end{table}

\begin{table}[H]
	\caption{Performance of SVM-classifier with pre-trained encoders for feature generation on the test dataset}
    \label{tab:test_perfomance_pre_svm}
    \begin{center}
        \begin{tabular}{c|c|c|c|c|c}
            \textbf{Feature Generator} & \textbf{Pooling Type} & \textbf{Accuracy} & \textbf{Precision} & \textbf{Recall} & \textbf{F1-score}  \\
            \hline
            CodeBERT & CLS & 0.24 & 0.66 & 0.24 & 0.16\\
            CodeBERT & Mean & 0.41 & 0.64 & 0.41 & 0.42\\
            UniXcoder & CLS & 0.41 & 0.67 & 0.41 & 0.46\\
            UniXcoder & Mean & 0.41 & 0.64 & 0.41 & 0.44\\
        \end{tabular}
    \end{center}
\end{table}


\begin{table}[H]
	\caption{Performance of SVM-classifier with fine-tuned encoders for feature generation on the test dataset}
    \label{tab:test_perfomance_fine_svm}
    \begin{center}
        \begin{tabular}{c|c|c|c|c|c}
            \textbf{Feature Generator} & \textbf{Pooling Type} & \textbf{Accuracy} & \textbf{Precision} & \textbf{Recall} & \textbf{F1-score}  \\
            \hline
            CodeBERT & CLS & 0.60 & 0.73 & 0.60 & 0.63\\
            CodeBERT & Mean & 0.64 & 0.71 & 0.64 & 0.66\\
            UniXcoder & CLS & 0.58 & 0.65 & 0.58 & 0.60\\
            UniXcoder & Mean & 0.57 & 0.65 & 0.57 & 0.60\\
        \end{tabular}
    \end{center}
\end{table}

\begin{table}[H]
    \caption{Performance of neural-network classifier on fine-trained encoders for feature generation on the test dataset}
    \label{tab:test_perfomance_fine_nn}
    \begin{center}
        \begin{tabular}{c|c|c|c|c|c}
            \textbf{Feature Generator} & \textbf{Pooling Type} & \textbf{Accuracy} & \textbf{Precision} & \textbf{Recall} & \textbf{F1-score}  \\
            \hline
            CodeBERT & CLS & 0.64 & 0.71 & 0.64 & 0.66\\
            CodeBERT & Mean & 0.66 & 0.72 & 0.66 & 0.68\\
            UniXcoder & CLS & 0.65 & 0.67 & 0.65 & 0.66\\
            UniXcoder & Mean & 0.62 & 0.68 & 0.62 & 0.64\\
        \end{tabular}
    \end{center}
\end{table}

\subsection{Experiment Discussion}

\noindent
In \crefrange{tab:test_perfomance_tfidf_svm}{tab:test_perfomance_fine_nn} we can see that our performance has dropped compared to the validation dataset, which was expected due to the bias in the validation dataset. The SVM classifier with tf-idf features achieved an F1-score of 0.51, representing the baseline performance we expect our models to reach. All SVM classifier models using pre-trained encoders fail to reach this baseline performance. However, this is expected, given the encoders have not been fine-tuned for this task. \\

\noindent
Pre-trained CodeBERT with CLS pooling performed extremely poorly. The model has an accuacy of $24\%$, which is worse than randomly guessing (i.e. $25\%$). This result indicates that without fine-tuning CodeBERT, CLS pooling produces useless embeddings for this task. The other configurations, however, can achieve an accuracy of $0.41$, far exceeding the performance of randomly guessing. This indicates that, while not performant enough for a real-world use case, these encoders can generate useful embeddings for this task without fine-tuning.\\

% \noindent
% Even more surprisingly is pre-trained UniXcoder's performance. In Section \ref{validation_discussion} it was predicted that pre-trained UnixCoder's high performance was a quirk of how SVM classifiers make predictions and would only be possible with a biased dataset. While this is likely true to some degree as pre-trained UniXcoder does not perform as well as the fine-tuned models on the test dataset (whereas in the validation dataset pre-trained UniXcoder and fine-tuned UniXcoder performed similarly), the model performs significantly better than randomly guessing and slightly better than an SVM classifier with tf-idf features. This is an interesting reseult. It would be worth investigating in future whether this is an isolated case; are there other classification tasks that UniXcoder (with a trained SVM classifier) can perform well in without fine-tuning? \\

% \noindent
% As with the validation dataset results, there is not a significant difference between the performance of the fine-tuned encoders with either an SVM or neural-network classification head. The neural-network classifier marginally outperforms the svm classifier models for all encoder and pooling configurations, except for CodeBERT with mean pooling. The best performing model was fine-tuned CodeBERT with mean pooling with a neural-network classification head with an F1-score of $0.68$. It is interesting that CodeBERT performs worse than UniXcoder without fine-tuning, but after fine-tuning is conducted CodeBERT slightly outperforms UniXcoder. \\

\noindent
As with the validation dataset results, there is not a significant difference between the performance of the fine-tuned encoders with either an SVM or neural network classifier. All neural-network classifier models marginally outperform all SVM classifier models, except for CodeBERT with mean pooling. The best performing model was CodeBERT with mean pooling using a neural-network classifier, which achieved an F1-score of $0.68$. This is interesting as UniXcoder performed better than CodeBERT before fine-tuning, but the opposite was true after fine-tuning. \\


\noindent
This experiment provides no strong evidence for a significant difference in performance between CLS and mean pooling when using fine-tuned encoders. While the difference was minimal, mean pooling performed better with CodeBERT, whereas CLS pooling performed better with UniXcoder. Recall from Section \ref{optimum_pooling} that the literature was inconclusive regarding whether one pooling strategy was generally better. 

\section{Performance On Unimodal Javadoc Models}
\label{javadoc_experiment}
\subsection{Experiment Description}

\noindent
So far, all model validation tests and experiments have been done on bimodal data. That is, both a Java method's Javadoc and code were included in the training data. Given this thesis' primary objective was to extract security-relevant features from \emph{Javadoc}, it is worth investigating how performant the models are when fine-tuned on Javadoc alone. \\

\noindent
For this experiment, the models will be evaluated on the test dataset once again, except the models will be fine-tuned on Javadoc only. \\

% \noindent
% In previous experiments, due to using bimodal data, if an input was longer than the maximum number of tokens into either codeBERT or uniXcoder ($512$ tokens), then both the code and natural language tokens would be truncated equally. For this experiment, we will not restrict the natural-language token length to $256$ but instead allow $512$ tokens as input. This means, for some inputs, the models will be trained on more Javadoc inputs than the previous bimodal experiments.

\subsection{Experiment Results}

\begin{table}[H]
    \caption{Performance of SVM-classifier with tf-idf feature generation on the test dataset (Javadoc only)}
    \label{tab:test_perfomance_tfidf_svm_javadoc}
        \begin{center}
            \begin{tabular}{c|c|c|c}
                \textbf{Accuracy} & \textbf{Precision} & \textbf{Recall} & \textbf{F1-score}  \\
                \hline
                0.59 & 0.52 & 0.59 & 0.55 \\
            \end{tabular}
        \end{center}
    \end{table}
    
    \begin{table}[H]
        \caption{Performance of SVM-classifier with pre-trained encoders for feature generation on the test dataset (Javadoc only)}
        \label{tab:test_perfomance_pre_svm_javadoc}
        \begin{center}
            \begin{tabular}{c|c|c|c|c|c}
                \textbf{Feature Generator} & \textbf{Pooling Type} & \textbf{Accuracy} & \textbf{Precision} & \textbf{Recall} & \textbf{F1-score}  \\
                \hline
                CodeBERT & CLS & 0.18 & 0.03 & 0.18 & 0.05 \\
                CodeBERT & Mean & 0.25 & 0.48 & 0.25  & 0.21 \\
                UniXcoder & CLS & 0.34 & 0.51 & 0.34 & 0.39 \\
                UniXcoder & Mean & 0.39 & 0.50 & 0.39 & 0.43 \\
            \end{tabular}
        \end{center}
    \end{table}
    
    
    \begin{table}[H]
        \caption{Performance of SVM-classifier with fine-tuned encoders for feature generation on the test dataset (Javadoc only)}
        \label{tab:test_perfomance_fine_svm_javadoc}
        \begin{center}
            \begin{tabular}{c|c|c|c|c|c}
                \textbf{Feature Generator} & \textbf{Pooling Type} & \textbf{Accuracy} & \textbf{Precision} & \textbf{Recall} & \textbf{F1-score}  \\
                \hline
                CodeBERT & CLS & 0.43 & 0.55 & 0.43 & 0.47 \\
                CodeBERT & Mean & 0.43 & 0.52 & 0.43 & 0.47 \\
                UniXcoder & CLS & 0.28 & 0.50 & 0.28 & 0.33 \\
                UniXcoder & Mean & 0.44 & 0.53 & 0.44 & 0.48 \\
            \end{tabular}
        \end{center}
    \end{table}
    
    \begin{table}[H]
        \caption{Performance of neural-network classifier on fine-trained encoders for feature generation on the test dataset (Javadoc Only)}
        \label{tab:test_perfomance_fine_nn_javadoc}
        \begin{center}
            \begin{tabular}{c|c|c|c|c|c}
                \textbf{Feature Generator} & \textbf{Pooling Type} & \textbf{Accuracy} & \textbf{Precision} & \textbf{Recall} & \textbf{F1-score}  \\
                \hline
                CodeBERT & CLS & 0.56 & 0.75 & 0.56 & 0.60 \\
                CodeBERT & Mean & 0.54 & 0.71 & 0.54 & 0.56 \\
                UniXcoder & CLS & 0.55 & 0.67 & 0.55 & 0.58 \\
                UniXcoder & Mean & 0.62 & 0.74 & 0.62 & 0.65 \\
            \end{tabular}
        \end{center}
\end{table}

\subsection{Experiment Discussion}

\noindent
It is clear from \crefrange{tab:test_perfomance_tfidf_svm_javadoc}{tab:test_perfomance_fine_nn_javadoc} that the performance of the unimodal models was worse than their bimodal equivalent. \\

\noindent
The SVM-Transformer models performed poorly. In particular, pre-trained CodeBERT with CLS pooling using an SVM classifier performed dramatically worse than randomly guessing with an accuracy of $0.18$ and an F1-score of $0.05$. Upon inspection of the predictions, it was found this model predicted all examples as \emph{sources}. It is unclear why the SVM classifier could not converge on a solution. However, it is clear that, without fine-tuning, CodeBERT cannot produce sufficient embeddings for this task. \\

\noindent
Another outlier was fine-tuned UniXcoder with CLS pooling using an SVM classifier. It performed significantly worse than other models in its category. The cause of this anomaly is unknown. \\

\noindent
The neural network classifier models (Table \ref{tab:test_perfomance_fine_nn_javadoc}) performed reasonably. They did not perform better than their bimodal equivalents (Table \ref{tab:test_perfomance_fine_nn}). However, this was to be expected given that these encoders were not pre-trained on Javadoc in isolation (i.e., without the corresponding code) \cite{CodeBERT}\cite{unixcoder}.\\

\noindent
Overall, since the results in Table \ref{tab:test_perfomance_fine_nn_javadoc} show performance far exceeding randomly guessing, it can be concluded that CodeBERT and UniXcoder \emph{are} learning and extracting security-relevant features from Javadoc. This is a significant result as it indicates the assumption made at the beginning of this thesis, that Javadoc contains security-relevant information, holds for our benchmarks.  

\section{Performance On Unimodal Code Models}

\subsection{Experiment Description}

\noindent
As was evident in Section \ref{javadoc_experiment}, security-relevant information can be extracted from Javadoc and generate useful results. To further investigate how much the code contributes to learning in the bimodal models, we conduct the test dataset experiment again, except we include only the method's code during fine-tuning and testing.

\subsection{Experiment Results}

\crefrange{tab:test_perfomance_tfidf_svm_code}{tab:test_perfomance_fine_nn_code} shows the performance of the model configurations on the test dataset when only code is used during fine-tuning and testing. \\

\noindent
In cases where the result is `N/A', this is where the model was not able to converge on a solution and generated NaN (not a number) embeddings. SVMs cannot classify NaNs, so there is no result. 

\begin{table}[H]
    \caption{Performance of SVM-classifier with tf-idf feature generation on the test dataset (Code only)}
    \label{tab:test_perfomance_tfidf_svm_code}
        \begin{center}
            \begin{tabular}{c|c|c|c}
                \textbf{Accuracy} & \textbf{Precision} & \textbf{Recall} & \textbf{F1-score}  \\
                \hline
                0.46 & 0.52 & 0.46 & 0.48 \\
            \end{tabular}
        \end{center}
    \end{table}
    
    \begin{table}[H]
        \caption{Performance of SVM-classifier with pre-trained encoders for feature generation on the test dataset (Code only)}
        \label{tab:test_perfomance_pre_svm_code}
        \begin{center}
            \begin{tabular}{c|c|c|c|c|c}
                \textbf{Feature Generator} & \textbf{Pooling Type} & \textbf{Accuracy} & \textbf{Precision} & \textbf{Recall} & \textbf{F1-score}  \\
                \hline
                CodeBERT & CLS & 0.18 & 0.03 & 0.18 & 0.05 \\
                CodeBERT & Mean & N/A & N/A & N/A & N/A\\
                UniXcoder & CLS & 0.29 & 0.5 & 0.29 & 0.35 \\
                UniXcoder & Mean & N/A & N/A & N/A & N/A\\
            \end{tabular}
        \end{center}
    \end{table}
    
    
    \begin{table}[H]
        \caption{Performance of SVM-classifier with fine-tuned encoders for feature generation on the test dataset (Code only)}
        \label{tab:test_perfomance_fine_svm_code}
        \begin{center}
            \begin{tabular}{c|c|c|c|c|c}
                \textbf{Feature Generator} & \textbf{Pooling Type} & \textbf{Accuracy} & \textbf{Precision} & \textbf{Recall} & \textbf{F1-score}  \\
                \hline
                CodeBERT & CLS & 0.49 & 0.52 & 0.49 & 0.50 \\
                CodeBERT & Mean & N/A & N/A & N/A & N/A\\
                UniXcoder & CLS & 0.41 & 0.48 & 0.41 & 0.44 \\
                UniXcoder & Mean & N/A & N/A & N/A & N/A\\
            \end{tabular}
        \end{center}
    \end{table}
    
    \begin{table}[H]
        \caption{Performance of neural-network classifier on fine-trained encoders for feature generation on the test dataset (Code only)}
        \label{tab:test_perfomance_fine_nn_code}
        \begin{center}
            \begin{tabular}{c|c|c|c|c|c}
                \textbf{Feature Generator} & \textbf{Pooling Type} & \textbf{Accuracy} & \textbf{Precision} & \textbf{Recall} & \textbf{F1-score}  \\
                \hline
                CodeBERT & CLS & 0.62 & 0.70 & 0.62 & 0.65 \\
                CodeBERT & Mean & 0.08 & 0.01 & 0.08 & 0.01 \\
                UniXcoder & CLS & 0.64 & 0.66 & 0.64 & 0.64 \\
                UniXcoder & Mean & 0.64 & 0.68 & 0.64 & 0.65 \\
            \end{tabular}
        \end{center}
\end{table}

\subsection{Experiment Discussion}

\noindent
The results of this experiment are interesting. In Table \ref{tab:test_perfomance_pre_svm_code} and Table \ref{tab:test_perfomance_fine_svm_code} both CodeBERT and Unixcoder with mean pooling were unable to converge on a solution. Furthermore, CodeBERT with mean pooling using a neural-network classifier had a performance significantly worse than randomly guessing (i.e. an accuracy of $0.08$).\\

\noindent
These results would indicate that mean pooling is not sufficient when operating on unimodal code data, however UniXcoder with mean pooling using a neural network classifier performs best in this experiment.\\

\noindent
For the neural network classifier models, where the model can converge on a solution, the models perform better than their equivalent in the Javadoc only experiment, i.e., Table \ref{tab:test_perfomance_fine_nn_javadoc}. \\

\noindent
Overall, these results are inconclusive. This experiment shows that while in some cases the Javadoc is not required to produce close to optimum results (compared to the highest performing bimodal models), for some model configurations, features extracted from code alone are useless (e.g., CodeBERT using mean pooling).

%\chapter{Results}
%Lorem ipsum dolor sit amet, consectetur adipiscing elit. Vivamus nulla libero, semper id dui finibus, suscipit pretium purus. Duis vehicula eleifend justo, at cursus turpis vulputate eget. Nullam vestibulum id neque a convallis. Sed molestie vehicula turpis, vel aliquet felis efficitur a. Fusce accumsan metus augue, id sagittis eros elementum vitae. Cras vel ultrices odio. Sed ac augue ex. Etiam tristique dignissim nibh, sit amet tincidunt libero tempor sit amet. Suspendisse tincidunt vitae ligula a posuere. Morbi euismod felis quis risus malesuada cursus. Curabitur imperdiet pretium odio, ac posuere turpis vehicula in. Praesent venenatis urna in augue congue elementum. Nam ligula nibh, rhoncus vitae luctus vel, feugiat sed neque. Curabitur ut arcu sed magna vulputate sollicitudin id a libero. Praesent tempus leo vel eros ornare mollis. Sed suscipit at mi ut viverra.


\chapter{Discussion}
\section{Performance}

\noindent
It was found that natural language processing on Javadoc is useful for classifying security-sensitive methods. A fine-tuned UniXcoder model was able to classify security-sensitive methods based on their Javadoc alone with an F1-score of $0.65$. \\

\noindent
However, it was also found that models trained on Javadocs exclusively did not perform as well as models using bimodal data (i.e., Javadoc \emph{and} code). This is likely a result of the specific pre-trained encoders being used. Since CodeBERT and UniXcoder were trained primarily on bimodal inputs, it is expected that these encoders perform best on input that more closely resembles the pre-training data. \\

\noindent
In almost all cases, the more sophisticated models (fine-tuned Transformer encoders with neural network classifiers) performed better than the more naive approaches (SVM classifiers, tf-idf feature generation, and pre-trained encoders without fine-tuning). In the case of models with bimodal input, the most simplistic model (an SVM classifier with tf-idf for feature extraction) had an F1-score of $0.51$ on the test dataset, whereas the best performing model (CodeBERT with mean pooling using a neural network classifier) achieved an F1-score of $0.68$. \\

\noindent
The performance of the Transformer encoders without fine-tuning was surprising. UniXcoder without fine-tuning, when paired with an SVM classifier, had an accuracy of $41\%$, which is considerably better than randomly guessing (i.e., $25\%$). In fact, the only model of this category that performed worse than randomly guessing was CodeBERT using CLS pooling which had an accuracy of $24\%$.\\

\noindent
Surprisingly, there were many cases of a significant difference in performance when using mean pooling versus CLS pooling. In general, it was found that UniXcoder performed better with mean pooling, and CodeBERT performed better with CLS pooling. At this stage, the cause of this discrepancy is unclear. \\

\noindent
Our models' performances fall short of Sas et al, SuSi, and SWAN, which achieved F1-scores of $86.4\%$, $91.9\%$, and $83\%$, respectively. However, this does not take away from the success of this thesis project. This thesis has shown that Javadoc contains information useful for security-related tasks, namely security-sensitive method classification. The approach presented in this thesis is not mutually exclusive to existing implementations and, rather than replace, could enhance current approaches. Furthermore, our models could be used by security analysts to generate a preliminary taint analysis configuration which could then be manually refined. This would likely result in a significant reduction in the time and effort required to produce a taint analysis configuration compared to creating one entirely from scratch. 

\section{Limitations and Improvements}
\label{limitations_and_improvements}

\noindent
There are some limitations to our approach. As is common with most machine learning projects, there was not enough high-quality labelled data to train and validate the models optimally. Ideally, the dataset should be larger and come from a greater diversity of projects. This would hopefully allow the models to generalise more effectively and perform better.\\

\noindent
Due to this lack of data, many examples needed to be manually labelled for this project. This is not ideal and undoubtedly lowered the quality of the dataset. A third party should review the dataset to ensure there are no mislabelled examples and that it is high-quality.\\

\noindent
Another limitation is in regards to how overloaded methods were handled. For this thesis, if a method was considered security-sensitive, all overloaded methods sharing its name were also considered security-sensitive. It is unclear if this assumption is always valid and is, therefore, a potential source of inaccuracy. \\

% \chapter{Milestone Task 1}
% Lorem ipsum dolor sit amet, consectetur adipiscing elit. Vivamus nulla libero, semper id dui finibus, suscipit pretium purus. Duis vehicula eleifend justo, at cursus turpis vulputate eget. Nullam vestibulum id neque a convallis. Sed molestie vehicula turpis, vel aliquet felis efficitur a. Fusce accumsan metus augue, id sagittis eros elementum vitae. Cras vel ultrices odio. Sed ac augue ex. Etiam tristique dignissim nibh, sit amet tincidunt libero tempor sit amet. Suspendisse tincidunt vitae ligula a posuere. Morbi euismod felis quis risus malesuada cursus. Curabitur imperdiet pretium odio, ac posuere turpis vehicula in. Praesent venenatis urna in augue congue elementum. Nam ligula nibh, rhoncus vitae luctus vel, feugiat sed neque. Curabitur ut arcu sed magna vulputate sollicitudin id a libero. Praesent tempus leo vel eros ornare mollis. Sed suscipit at mi ut viverra.

Maecenas vitae tempus arcu. Nullam maximus id ligula a finibus. Aliquam laoreet faucibus sem. Mauris interdum erat vel pretium iaculis. Sed et pretium arcu, at condimentum lorem. Suspendisse tincidunt neque vel ullamcorper consequat. Phasellus eget ligula quis velit fringilla porttitor. Pellentesque quis risus id erat euismod bibendum mattis nec sapien. Duis gravida, turpis at dapibus consequat, ex mi tempor augue, in scelerisque mi nisi ut lorem. Donec orci tellus, lobortis ac pretium non, elementum et dolor. Morbi sed hendrerit erat. Praesent nec ante lorem. Vestibulum dapibus varius nisl at hendrerit. Praesent aliquam vehicula dolor et pulvinar. Aenean facilisis leo eget lacus rhoncus, ac tempus urna molestie. Vivamus efficitur erat quis felis pulvinar mattis quis volutpat lorem.

Nullam luctus velit varius massa suscipit fringilla. Quisque egestas diam nec eros sagittis posuere. Sed venenatis facilisis nibh ac mattis. Morbi a mi at est vehicula imperdiet. Fusce tempus justo at tortor volutpat vulputate. Vestibulum molestie congue interdum. Nullam auctor faucibus risus ac hendrerit. Nulla viverra fermentum vulputate. Nulla mattis elementum ligula. Fusce eu tellus sapien. Maecenas scelerisque odio lacus, quis venenatis sapien molestie semper. Donec in consequat libero, quis euismod tellus.

Vestibulum eget lobortis lorem, id aliquam leo. Morbi ornare, enim et iaculis luctus, erat dui rutrum ligula, vel aliquam libero nisi vitae eros. Proin sodales ipsum nec ipsum tristique, porttitor sagittis tortor rhoncus. Integer auctor rutrum imperdiet. Nulla sollicitudin facilisis mauris, in mattis lorem ultrices quis. Curabitur neque lectus, auctor at ullamcorper et, porta eu nulla. Cras efficitur felis in ipsum volutpat, vitae condimentum velit ornare.

Donec consectetur, libero tempus sollicitudin gravida, felis erat aliquam libero, viverra maximus lacus nisi in lacus. Aenean hendrerit volutpat felis, sed viverra ex accumsan a. Vivamus vel tortor viverra nisi tincidunt dapibus. Proin diam libero, varius at turpis et, tincidunt congue arcu. Nunc tristique elit libero, nec rutrum quam aliquet at. Orci varius natoque penatibus et magnis dis parturient montes, nascetur ridiculus mus. Vestibulum at nisi quis nulla malesuada tempor. Mauris eu felis risus. Vestibulum ante ipsum primis in faucibus orci luctus et ultrices posuere cubilia curae; Donec convallis est ex, vitae facilisis eros commodo in. Nullam faucibus nunc ac magna congue, et porttitor magna eleifend. Nunc at velit ullamcorper, pellentesque lorem in, dictum arcu. Suspendisse nec auctor velit. Mauris congue magna molestie augue maximus, ac hendrerit ante viverra. Pellentesque convallis nisl et felis venenatis varius.

% \chapter{Milestone Task 2}
% Lorem ipsum dolor sit amet, consectetur adipiscing elit. Vivamus nulla libero, semper id dui finibus, suscipit pretium purus. Duis vehicula eleifend justo, at cursus turpis vulputate eget. Nullam vestibulum id neque a convallis. Sed molestie vehicula turpis, vel aliquet felis efficitur a. Fusce accumsan metus augue, id sagittis eros elementum vitae. Cras vel ultrices odio. Sed ac augue ex. Etiam tristique dignissim nibh, sit amet tincidunt libero tempor sit amet. Suspendisse tincidunt vitae ligula a posuere. Morbi euismod felis quis risus malesuada cursus. Curabitur imperdiet pretium odio, ac posuere turpis vehicula in. Praesent venenatis urna in augue congue elementum. Nam ligula nibh, rhoncus vitae luctus vel, feugiat sed neque. Curabitur ut arcu sed magna vulputate sollicitudin id a libero. Praesent tempus leo vel eros ornare mollis. Sed suscipit at mi ut viverra.

Maecenas vitae tempus arcu. Nullam maximus id ligula a finibus. Aliquam laoreet faucibus sem. Mauris interdum erat vel pretium iaculis. Sed et pretium arcu, at condimentum lorem. Suspendisse tincidunt neque vel ullamcorper consequat. Phasellus eget ligula quis velit fringilla porttitor. Pellentesque quis risus id erat euismod bibendum mattis nec sapien. Duis gravida, turpis at dapibus consequat, ex mi tempor augue, in scelerisque mi nisi ut lorem. Donec orci tellus, lobortis ac pretium non, elementum et dolor. Morbi sed hendrerit erat. Praesent nec ante lorem. Vestibulum dapibus varius nisl at hendrerit. Praesent aliquam vehicula dolor et pulvinar. Aenean facilisis leo eget lacus rhoncus, ac tempus urna molestie. Vivamus efficitur erat quis felis pulvinar mattis quis volutpat lorem.

Nullam luctus velit varius massa suscipit fringilla. Quisque egestas diam nec eros sagittis posuere. Sed venenatis facilisis nibh ac mattis. Morbi a mi at est vehicula imperdiet. Fusce tempus justo at tortor volutpat vulputate. Vestibulum molestie congue interdum. Nullam auctor faucibus risus ac hendrerit. Nulla viverra fermentum vulputate. Nulla mattis elementum ligula. Fusce eu tellus sapien. Maecenas scelerisque odio lacus, quis venenatis sapien molestie semper. Donec in consequat libero, quis euismod tellus.

Vestibulum eget lobortis lorem, id aliquam leo. Morbi ornare, enim et iaculis luctus, erat dui rutrum ligula, vel aliquam libero nisi vitae eros. Proin sodales ipsum nec ipsum tristique, porttitor sagittis tortor rhoncus. Integer auctor rutrum imperdiet. Nulla sollicitudin facilisis mauris, in mattis lorem ultrices quis. Curabitur neque lectus, auctor at ullamcorper et, porta eu nulla. Cras efficitur felis in ipsum volutpat, vitae condimentum velit ornare.

Donec consectetur, libero tempus sollicitudin gravida, felis erat aliquam libero, viverra maximus lacus nisi in lacus. Aenean hendrerit volutpat felis, sed viverra ex accumsan a. Vivamus vel tortor viverra nisi tincidunt dapibus. Proin diam libero, varius at turpis et, tincidunt congue arcu. Nunc tristique elit libero, nec rutrum quam aliquet at. Orci varius natoque penatibus et magnis dis parturient montes, nascetur ridiculus mus. Vestibulum at nisi quis nulla malesuada tempor. Mauris eu felis risus. Vestibulum ante ipsum primis in faucibus orci luctus et ultrices posuere cubilia curae; Donec convallis est ex, vitae facilisis eros commodo in. Nullam faucibus nunc ac magna congue, et porttitor magna eleifend. Nunc at velit ullamcorper, pellentesque lorem in, dictum arcu. Suspendisse nec auctor velit. Mauris congue magna molestie augue maximus, ac hendrerit ante viverra. Pellentesque convallis nisl et felis venenatis varius.

% obviously main experiments go here, later

\chapter{Conclusion}
% \noindent
% The goal of this thesis was to use machine-learning techniques to automatically classify security-sensitive Java methods using Javadoc. \\

% \noindent
% A dataset of security-relevant Java methods' source-code and Javadoc was produced. To our knowledge, a similar dataset does not already exist in the literature. The final dataset (training and validation combined) is comprised of 1096 methods (158 sources, 268 sinks, 175 sanitisers, and 495 security-insensitive methods). The dataset is produced dynamically from a preliminary dataset format. The preliminary dataset and the code to dynamically build the dataset is publicly available on \href{https://github.com/uqcyber/SecurityMethodClassifier/tree/main/tools/dataset-manager}{GitHub}. \\

% \noindent
% In addition to the training and validation datasets, a test dataset was created on a more realistic use-case. This test dataset contained a random sample of methods from Apache's Log4j ($n = 165$) which were manually classified. This test dataset was used to evaluate the models' performance.  \\

% \noindent
% The pre-trained models CodeBERT and UniXcoder were fine-tuned on our dataset and used to perform the classification. Both mean and CLS pooling strategies were used as it was not clear which was more effective for this use-case. It was found that CLS pooling performed better for UniXcoder and mean pooling performed better for CodeBERT. \\

% \noindent
% The highest performing model configuration was found to be fine-tuned CodeBERT with a neural network classifier using mean pooling. This model achieved an F1-score of $0.68$. This result is significantly better than the performance of a naive SVM classifier using tf-idf for feature generation, which had an F1-score of 0.51. \\

% % \noindent
% % To our knowledge, the approach we present in this thesis is novel. Furthermore, the results of this thesis indicate that Javadoc contains information useful for security analyses.

% \noindent
% The approach we have taken in this thesis, to our knowledge, is novel. While the models we have presented in this thesis are not accurate enough to allow for taint analysis configurations to be fully automatically generated, as was hoped at beginning of this thesis, the models do perform well enough to allow for security analysts to create configurations more efficiently than from scratch. Furthermore, it is possible our approach could be combined with existing approaches in the literature to produce an even more performant model. Most importantly, however, this thesis has shown that Javadoc does indeed contain information useful for security anaylsis and it is worth exploring what other security analyses or other tasks could be enhanced using deep learning techniques on Javadoc.

\noindent
In this concluding chapter of the thesis, we summarise the project's outcomes and discuss potential future work for this project.

\section{Project Outcomes}

\noindent
A dataset of security-relevant Java methods' source code and Javadoc was produced. To our knowledge, a similar dataset does not already exist in the literature. The final dataset (training and validation combined) is comprised of 1096 methods (158 sources, 268 sinks, 175 sanitisers, and 495 security-insensitive methods). The dataset is produced dynamically from a preliminary dataset. The preliminary dataset and the code to dynamically build the dataset is publicly available on \href{https://github.com/uqcyber/SecurityMethodClassifier/tree/main/tools/dataset-manager}{GitHub}. \\

\noindent
In addition to this dataset, a test dataset was created. This test dataset consisted of a random sample of Java methods from Apache's Log4J that were manually labelled. This test dataset served as the benchmark test for our models as it more accurately reflects a realistic use case for the models. \\

\noindent
Furthermore, several machine learning models based on CodeBERT and UniXcoder were produced. These pre-trained models were evaluated and tested in a variety of configurations. The difference between different pooling strategies --- mean versus CLS pooling --- was investigated. It was determined that UniXcoder performed best with CLS pooling, and CodeBERT performed best with mean pooling. \\

\noindent
The highest-performing model produced for this thesis project was CodeBERT with a neural-network classifier using mean pooling. This model achieved an F1-score of $0.68$ on the test dataset. The source code for all the machine-learning models and configurations produced for this thesis can be found on \href{https://github.com/uqcyber/SecurityMethodClassifier/tree/main/models}{GitHub}. \\


\section{Future Work}

\noindent
In addition to the improvements suggested in Section \ref{limitations_and_improvements}, there are a number of extensions to this project that can be undertaken. \\

\noindent
One possible extension is to include the \emph{context} of a method call into the training data. For instance, it may be useful to include the Javadoc for the method that calls the method to be classified into the training data. This is likely to be a difficult extension to implement as it would require a call-graph analysis to be conducted. \\ 

\noindent
Another potential extension would be to include in-line comments in the training data. So far, in-line comments have been removed, but perhaps further performance could be gained by including all the in-line comments of a method (in addition to the Javadoc) in the training data. Recall that the dataset provided for this thesis is not static; it is produced dynamically by the dataset manager. As such, implementing this extension would likely only require small tweaks to the dataset manager and for the models to be trained once again. \\

\noindent
Another simple extension is to `highlight' keywords in the input during training and ensure the tokeniser does not split them. Furthermore, these tokens could have their embeddings fixed so their embeddings are not altered by contextual information. It is not clear at this stage what the keywords would be, but it is worth exploring this idea in future.\\

\noindent
Finally, learning could be conducted on an additional classification task where the model learns to classify each token in the input as either security-sensitive or security-insensitive. This approach was inspired by Wang et. al's identifier-tagging approach \cite{identifier_tagging} and may allow for more sophisticated learning to occur.

\cleardoublepage\pagenumbering{Roman}
\appendix
%\chapter{Appendix}
\chapter{Default Taint Analysis Configuration}
\label{appendix:sources}


Figure \ref{sink-tank-default-source-config}, below, is a real default configuration of sources for the taint-analyis tool `sink-tank'. Importantly, from the example shown in Figure \ref{figure-taint-example}, we can see that HttpServletRequest.getParameter() was considered a source. Luckily, since ServletRequest.getParameter() is listed in the configuration, (recalling that HttpServletRequest implements the interface ServletRequest), it is likely the vulnerability would be detected by sink-tank.\\

\begin{figure}[H]
\lstinputlisting[language=XML]{misc/method-sources.xml}
\caption{A portion of the default configuration of sources for the taint-analysis tool \emph{sink-tank}}
\label{sink-tank-default-source-config}
\end{figure}

\newpage


\printbibliography

\end{document}